% CS301 Week 6 -- Singly linked lists: array-vs-linked comparison, node
% anatomy with structures, types of linked lists, applications, and the five
% core singly-linked operations (insert, delete, search, display, reverse).
% Difficulty: intro (course-wide default per knowledge-base-CS301.md).
% Page-grounded in Karumanchi2017 Ch.3 (Linked Lists, pp.74-162 PDF) and
% Sedgewick2011 p.142 ("Linked lists" definition); Wirth2004 for the
% pointer/dynamic-allocation concepts behind node creation; Horowitz & Sahni
% Ch.4 / Aho-Hopcroft-Ullman Ch.2 cited at chapter level only (not
% page-indexed).
% Handoff: Week 5 closed with the queue contract, the circular-buffer fix
% ((rear+1) mod n with the count convention, Lab P5), deques, array-based
% priority queues (Lab P6), and the Unit-I review for Class Test 1 -- and
% promised "what changes when nodes carry their own links instead of sitting
% in numbered slots". This deck does NOT re-teach queues or circular
% arithmetic -- it builds on them. Continues the expression-processor thread:
% the calculator can buffer jobs, but its identifier table still needs a home.
% Sets up Week 7 (doubly/circular lists, list-based stack/queue, polynomials).

\documentclass[aspectratio=169]{beamer}
\input{header}
\coursecode{CS301}

\title{Chains Instead of Slots}
\subtitle{Week 6 --- Singly Linked Lists: Nodes, Links, and Pointer Surgery}
\author[Dr. Auqib Hamid Lone]{\textbf{Dr. Auqib Hamid Lone}}
\institute[GCET Ganderbal]{PCC CS-301: Data Structures\\GCET, Ganderbal}
\date[Week 6]{\small Week 6}

\begin{document}

\frame{\titlepage}

% =============================================================================
% ACT 0 --- MOTIVATION + HANDOFF
% =============================================================================
\section{Motivation and the Handoff}

\begin{frame}{Today's Four Questions}
  \begin{enumerate}
    \item Why does inserting one name into a full array cost $n$ moves?
    \item What does a node carry that an array slot does not?
    \item Which linked shape fits which job --- singly, doubly, or circular?
    \item How do three pointers turn a chain around without copying it?
  \end{enumerate}
  \vspace{0.2cm}
  \begin{exampleblock}{Plan}
    Price the array's weakness. Build the node. Tour the family. Operate:
    insert, delete, search, display --- then reverse.
  \end{exampleblock}
\end{frame}

\begin{frame}{Handoff: What Week 5 Finished}
  \begin{itemize}
    \item Week 5 gave us the queue contract (\texttt{enqueue} at
      \texttt{rear}, \texttt{dequeue} from \texttt{front}), the circular fix
      --- \texttt{rear} $=$ (\texttt{rear} $+ 1) \bmod n$ with the
      \texttt{count} convention (Lab P5) --- deques, and array-based priority
      queues (Lab P6), then closed Unit I for Class Test 1.
    \item Every structure so far lives in \emph{numbered slots}: an array owns
      its elements by position, and position is both its power and its tax.
    \item Its closing promise: ``what changes when nodes carry their own
      links instead of sitting in numbered slots.''
  \end{itemize}
  \begin{exampleblock}{This week's job}
    Free the elements from their slots: each one carries the address of the
    next. Price what we gain (cheap insert/delete) against what we lose
    (random access).
  \end{exampleblock}
\end{frame}

\begin{frame}{The Calculator Needs a Symbol Table}
  \begin{itemize}
    \item Weeks 3--5 built an expression processor: tokenize, check brackets,
      convert, evaluate, buffer jobs in arrival order. One program learning
      new tricks each week.
    \item New deficiency: evaluating \texttt{(a+b)*c} needs the \emph{values}
      of \texttt{a}, \texttt{b}, \texttt{c} --- identifiers arrive at
      unpredictable times, in unpredictable numbers. A fixed array demands we
      guess the capacity up front.
    \item The same need appears everywhere: OS free-lists of memory blocks,
      browser back-forward chains, music playlists, and (Week 7) polynomials
      with missing terms.
  \end{itemize}
  \begin{exampleblock}{The pattern (design principle)}
    Every structure in this course answers a deficiency the previous one just
    exhibited. The array's deficiency is fixed capacity plus shifting --- the
    linked list is the answer.
  \end{exampleblock}
\end{frame}

\begin{frame}{Where This Week Sits}
  \begin{itemize}
    \item Week 1 gave us \texttt{malloc}/\texttt{free} and \texttt{NULL};
      Week 2 the pricing habit (operation $+$ bound $+$ case). Every claim
      today repeats that shape.
    \item Week 7 re-implements the stack and the queue on top of \emph{this}
      week's chain, adds the backward link (\texttt{p->prev}), and spends the
      polynomial application; Week 10 replays the same symbol table as a hash
      table --- $O(n)$ as a list vs.\ $O(1)$ expected hashed.
    \item New symbols today: \texttt{head} (first-node pointer,
      \texttt{NULL} when empty) and \texttt{p->next} (the successor link).
  \end{itemize}
  \muted{The arc again: state the contract, implement it once, analyse what
    it costs.}
\end{frame}

\transitionslide{First: price the slot tax that makes chains worth building.}

% =============================================================================
% ACT 1 --- ARRAY VS LINKED
% =============================================================================
\section{Why Chains}

\sectiondivider{Section 1}{Why Chains}

\begin{frame}{The Slot Tax: Insert Into a Full Array}
  \begin{itemize}
    \item An array of $n$ integers holds \texttt{A[0..n-1]} back to back.
      Insert \texttt{x} at position $k$ and every element from $k$ to $n-1$
      must shift one slot right --- up to $n$ moves, $O(n)$ worst case.
    \item Delete at $k$ and the same crowd shifts left. Growing past capacity
      means \texttt{malloc} a bigger block, copy all $n$, \texttt{free} the
      old one.
    \item \bad{Position is the tax}: the array finds any element in $O(1)$
      worst case by arithmetic, but pays linear time whenever the lineup
      itself changes.
  \end{itemize}
  \begin{exampleblock}{The question that buys the list}
    What if each element carried the address of its successor, so the lineup
    could change by rewriting one address instead of moving $n$ elements?
  \end{exampleblock}
\end{frame}

\begin{frame}{Array Slots vs.\ Linked Nodes, Drawn}
  \begin{center}
    \begin{tikzpicture}[every node/.style={font=\scriptsize}]
      \node[text=neutral] at (-4.6, 1.15) {array \texttt{A[0..3]}};
      \node[draw, fill=light-bg, minimum width=1.1cm, minimum height=0.7cm,
        text width=0.9cm, align=center] (a0) at (-4.6, 0.45) {10};
      \node[draw, fill=light-bg, minimum width=1.1cm, minimum height=0.7cm,
        text width=0.9cm, align=center] (a1) at (-3.3, 0.45) {20};
      \node[draw, fill=light-bg, minimum width=1.1cm, minimum height=0.7cm,
        text width=0.9cm, align=center] (a2) at (-2.0, 0.45) {30};
      \node[draw, fill=light-bg, minimum width=1.1cm, minimum height=0.7cm,
        text width=0.9cm, align=center] (a3) at (-0.7, 0.45) {40};
      \node[text=neutral, below=0.15cm of a1] {neighbours by arithmetic};

      \node[text=neutral] at (3.0, 1.15) {list: neighbours by address};
      \node[draw, fill=white, minimum width=0.9cm, minimum height=0.7cm,
        align=center] (n10) at (1.3, -0.35) {10};
      \node[draw, fill=white, minimum width=0.9cm, minimum height=0.7cm,
        align=center] (n11) at (2.2, -0.35) {$\to$};
      \node[draw, fill=white, minimum width=0.9cm, minimum height=0.7cm,
        align=center] (n20) at (3.3, -0.35) {20};
      \node[draw, fill=white, minimum width=0.9cm, minimum height=0.7cm,
        align=center] (n21) at (4.2, -0.35) {$\to$};
      \node[draw, fill=white, minimum width=0.9cm, minimum height=0.7cm,
        align=center] (n30) at (5.3, -0.35) {30};
      \node[draw, fill=neutral!20, minimum width=0.9cm, minimum height=0.7cm,
        align=center] (n31) at (6.2, -0.35) {\texttt{NULL}};
      \draw[->, thick, color=primary-blue] (0.2, -0.35) -- (n10.west);
      \node[text=primary-blue] at (0.2, 0.15) {\texttt{head}};
      \node[text=neutral, below=0.15cm of n20] {each box lives anywhere};
    \end{tikzpicture}
  \end{center}
  \vspace{-0.2cm}
  \begin{itemize}
    \item Top: neighbours sit side by side, found by index arithmetic.
      Bottom: neighbours sit \emph{anywhere} in the heap, found by following
      arrows from \texttt{head} to \texttt{NULL}.
  \end{itemize}
\end{frame}

\begin{frame}{The Node: Data Plus One Address}
  \begin{block}{Definition (Sedgewick Sec.~1.3, p.~142; Karumanchi Ch.~3,
    pp.74--162 PDF)
    \cite{Sedgewick2011_algorithms,Karumanchi2017_data_structures_algorithms_made_easy}}
    A \key{linked list} is either empty (\texttt{head} $=$ \texttt{NULL}) or
    a node holding data plus a reference to the rest of the list. In C:
  \end{block}
\begin{center}
{\small \texttt{struct node \{ int data; struct node *next; \};}}
\end{center}
  \begin{itemize}
    \item Each node is separately \texttt{malloc}'d --- Week 1's allocation
      trace, one box per element, scattered across the heap (stack grows
      down, heap up, same orientation as Week 1).
    \item \texttt{head} is the \emph{only} entry point: lose it and the whole
      chain leaks. \texttt{p->next} of the last node is \texttt{NULL}.
  \end{itemize}
\end{frame}

\begin{frame}{Worked Example: Build 10--20--30 by Hand}
  \begin{center}
  {\small
  \begin{tabular}{ll}
    \toprule
    \textbf{step} & \textbf{heap after the step} \\
    \midrule
    \texttt{head = NULL} & empty: \texttt{head} points nowhere \\
    insert 10 at front & \texttt{head} $\to$ \texttt{[10|NULL]} \\
    insert 20 at front & \texttt{head} $\to$ \texttt{[20|$\cdot$]} $\to$
      \texttt{[10|NULL]} \\
    insert 30 at end & \texttt{head} $\to$ 20 $\to$ 10 $\to$ 30 $\to$
      \texttt{NULL} \\
    \bottomrule
  \end{tabular}}
  \end{center}
  \begin{itemize}
    \item Read the arrows, not the addresses: the three boxes may sit at any
      heap locations --- order lives in the links, never in the layout.
    \item Every arrow written is one pointer assignment; every box is one
      \texttt{malloc} that a matching \texttt{free} must one day reclaim.
  \end{itemize}
\end{frame}

\begin{frame}{Arrays vs.\ Linked Lists: The Honest Bill}
  \begin{center}
  {\footnotesize
  \begin{tabular}{lll}
    \toprule
    \textbf{operation} & \textbf{array} & \textbf{singly linked list} \\
    \midrule
    access $i$-th element & $O(1)$ worst case & $O(n)$ worst case (walk $i$ links) \\
    insert/delete at front & $O(n)$ worst case (shift) & $O(1)$ worst case \\
    insert/delete at end & $O(1)$* amortized & $O(n)$ worst case (no tail link) \\
    memory & fixed, contiguous & grows one node at a time \\
    \bottomrule
  \end{tabular}}
  \end{center}
  \begin{itemize}
    \item *Append is $O(1)$ amortized only with spare capacity; a full array
      pays the copy-everything reallocation.
    \item No winner --- only workloads: many positional reads $\to$ array;
      many front insertions and deletions $\to$ list. State the workload,
      then choose.
  \end{itemize}
\end{frame}

\transitionslide{The chain is priced. Now meet the whole family.}

% =============================================================================
% ACT 2 --- TYPES + APPLICATIONS
% =============================================================================
\section{Types and Applications}

\sectiondivider{Section 2}{Types and Applications}

\begin{frame}{Three Shapes of Chain}
  \begin{center}
    \begin{tikzpicture}[every node/.style={font=\scriptsize}]
      \node[text=neutral] at (-4.5, 1.3) {singly};
      \node[draw, fill=white, minimum width=0.8cm, minimum height=0.6cm,
        align=center] (s1) at (-5.3, 0.55) {a};
      \node[draw, fill=white, minimum width=0.8cm, minimum height=0.6cm,
        align=center] (s2) at (-4.3, 0.55) {b};
      \node[draw, fill=white, minimum width=0.8cm, minimum height=0.6cm,
        align=center] (s3) at (-3.3, 0.55) {c};
      \draw[->, thick, color=jet] (s1.east) -- (s2.west);
      \draw[->, thick, color=jet] (s2.east) -- (s3.west);
      \draw[->, thick, color=neutral] (s3.east) -- (s3.east -| 2.2, 0.55);
      \node[text=neutral, below=0.15cm of s2] {one way; \texttt{NULL} ends it};

      \node[text=neutral] at (0.3, 1.3) {doubly (Week 7)};
      \node[draw, fill=white, minimum width=0.8cm, minimum height=0.6cm,
        align=center] (d1) at (-0.5, 0.55) {a};
      \node[draw, fill=white, minimum width=0.8cm, minimum height=0.6cm,
        align=center] (d2) at (0.5, 0.55) {b};
      \node[draw, fill=white, minimum width=0.8cm, minimum height=0.6cm,
        align=center] (d3) at (1.5, 0.55) {c};
      \draw[<->, thick, color=jet] (d1.east) -- (d2.west);
      \draw[<->, thick, color=jet] (d2.east) -- (d3.west);
      \node[text=neutral, below=0.15cm of d2] {\texttt{prev} + \texttt{next}};

      \node[text=neutral] at (4.3, 1.3) {circular (Week 7)};
      \node[draw, fill=white, minimum width=0.8cm, minimum height=0.6cm,
        align=center] (c1) at (3.6, 0.55) {a};
      \node[draw, fill=white, minimum width=0.8cm, minimum height=0.6cm,
        align=center] (c2) at (4.6, 0.55) {b};
      \node[draw, fill=white, minimum width=0.8cm, minimum height=0.6cm,
        align=center] (c3) at (5.6, 0.55) {c};
      \draw[->, thick, color=jet] (c1.east) -- (c2.west);
      \draw[->, thick, color=jet] (c2.east) -- (c3.west);
      \draw[->, thick, color=jet] (c3.south) -- (4.6, -0.35) -- (c1.south);
      \node[text=neutral, below=0.55cm of c2] {last points home; no \texttt{NULL}};
    \end{tikzpicture}
  \end{center}
  \vspace{-0.2cm}
  \begin{itemize}
    \item \key{Singly} (today): one link each, cheapest per node.
      \key{Doubly}: backward walks, twice the link maintenance.
      \key{Circular}: round-robin buffers and schedulers.
  \end{itemize}
\end{frame}

\begin{frame}{Applications: Where Chains Earn Their Keep}
  \begin{block}{The rule of thumb (Karumanchi Ch.~3, pp.74--162 PDF; standard
    treatment)
    \cite{Karumanchi2017_data_structures_algorithms_made_easy}}
    Reach for a list when the \key{number of items is unpredictable} and
    \key{insertions and deletions land anywhere} --- the two places arrays
    charge the slot tax.
  \end{block}
  \begin{itemize}
    \item \key{Symbol table} (our calculator): identifiers come and go;
      lookup walks the chain --- $O(n)$ worst case, the bill Week 10's hash
      table will cut to $O(1)$ expected.
    \item \key{OS memory management}: free blocks linked as they are freed
      and re-split on demand --- no fixed table of blocks.
    \item \key{Polynomials} (Week 7 preview): $5x^{100} + 3x^{2} + 1$ stores
      three nodes, not 101 array slots mostly holding zero.
  \end{itemize}
\end{frame}

\begin{frame}{Check Your Prediction}
  \begin{block}{A text editor holds 1\,000 lines. The user inserts one line at
    the top, 500 times. Array or singly linked list?}
    A: array (indexing wins) \quad B: linked list \quad
    C: identical either way
  \end{block}
  \begin{enumerate}
    \item A --- indexing only matters if we jump to line numbers; here the
      workload is all front insertions.
    \item B --- each insert rewrites one \texttt{next} pointer: $O(1)$ worst
      case, where the array shifts up to $n$ lines every time.
    \item C --- the table two sections back says otherwise; price the
      workload, not the structure.
  \end{enumerate}
  \muted{Correct answer: B. Reads are sequential either way; the insertions
    decide.}
\end{frame}

\transitionslide{The family introduced. Now the surgery: five operations, one link at a time.}

% =============================================================================
% ACT 3 --- CORE OPERATIONS
% =============================================================================
\section{Insert, Delete, Search, Display}

\sectiondivider{Section 3}{Insert, Delete, Search, Display}

\begin{frame}[fragile]{Insert at the Front: Two Assignments}
  \begin{block}{Pointer surgery (Karumanchi Ch.~3, pp.74--162 PDF)}
\begin{semiverbatim}
newNode = malloc(sizeof(struct node))  // 1. fresh box
newNode->data = x                      // 2. fill it
newNode->next = head                   // 3. point at old chain FIRST
head = newNode                         // 4. THEN swing head over
\end{semiverbatim}
  \end{block}
  \begin{itemize}
    \item Order is the whole lesson: step 4 before step 3 orphans the chain
      --- \texttt{head}'s old value is gone with no copy anywhere.
    \item Cost: two pointer writes, no walk --- $O(1)$ worst case. Always
      check \texttt{malloc} for \texttt{NULL} before dereferencing.
  \end{itemize}
\end{frame}

\begin{frame}{Insert at the Front, Drawn}
  \begin{center}
    \begin{tikzpicture}[every node/.style={font=\scriptsize}]
      \node[draw, fill=primary-gold!25, minimum width=1.0cm,
        minimum height=0.7cm, align=center] (nw) at (0, 0) {x};
      \node[draw, fill=white, minimum width=1.0cm, minimum height=0.7cm,
        align=center] (o1) at (3.0, 0) {10};
      \node[draw, fill=white, minimum width=1.0cm, minimum height=0.7cm,
        align=center] (o2) at (4.5, 0) {20};
      \node[draw, fill=neutral!20, minimum width=1.0cm, minimum height=0.7cm,
        align=center] (o3) at (6.0, 0) {\texttt{NULL}};
      \draw[->, thick, color=jet] (o1.east) -- (o2.west);
      \draw[->, thick, color=jet] (o2.east) -- (o3.west);

      \draw[->, thick, color=positive] (nw.east) -- (1.4, 0.55) -- (o1.north);
      \node[text=positive, above=0.15cm of nw] {step 3 first};

      \draw[->, thick, dashed, color=primary-blue] (-1.2, -0.9) -- (nw.south);
      \node[text=primary-blue, anchor=east] at (-1.2, -0.9) {\texttt{head}, step 4};

      \node[text=negative, below=0.55cm of o1]
        {swap the steps and the old chain is unreachable};
    \end{tikzpicture}
  \end{center}
  \vspace{-0.2cm}
  \begin{itemize}
    \item Trace it with a finger: new box points at the old first node,
      \emph{then} \texttt{head} moves. Either arrow alone is a broken list.
  \end{itemize}
\end{frame}

\begin{frame}[fragile]{Insert at a Position and at the End}
  \begin{block}{The walk-then-splice pattern}
\begin{semiverbatim}
p = head; for (i = 0; i < pos-1; i++) p = p->next
newNode->next = p->next; p->next = newNode
\end{semiverbatim}
  \end{block}
  \begin{itemize}
    \item Walk to the predecessor ($O(n)$ worst case), then the same two
      writes as front-insert in the same order.
    \item End-insert is position $n$: walk all $n$ links --- $O(n)$ worst
      case without a tail pointer. Guard: inserting into an empty list sets
      \texttt{head} directly; walking past \texttt{NULL} is a bug, not an
      edge case.
  \end{itemize}
\end{frame}

\begin{frame}[fragile]{Delete: Find the Predecessor, Bypass, Free}
  \begin{block}{Delete first occurrence of \texttt{x}}
\begin{semiverbatim}
if head->data == x: temp = head; head = head->next
else: find p with p->next->data == x
      temp = p->next; p->next = temp->next
free(temp)
\end{semiverbatim}
  \end{block}
  \begin{itemize}
    \item Three acts: \key{bypass} (rewrite one link), \key{reclaim}
      (\texttt{free} exactly once), \key{never touch} \texttt{temp} after the
      free --- that is Week 1's dangling pointer by name.
    \item Search-then-delete costs the search: $O(n)$ worst case; the bypass
      itself is $O(1)$ worst case once \texttt{p} is in hand.
  \end{itemize}
\end{frame}

\begin{frame}{Delete the Middle, Drawn}
  \begin{center}
    \begin{tikzpicture}[every node/.style={font=\scriptsize}]
      \node[draw, fill=white, minimum width=1.0cm, minimum height=0.7cm,
        align=center] (m1) at (0, 0) {10};
      \node[draw, fill=negative!15, minimum width=1.0cm, minimum height=0.7cm,
        align=center] (m2) at (2.0, 0) {20};
      \node[draw, fill=white, minimum width=1.0cm, minimum height=0.7cm,
        align=center] (m3) at (4.0, 0) {30};
      \node[draw, fill=neutral!20, minimum width=1.0cm, minimum height=0.7cm,
        align=center] (m4) at (6.0, 0) {\texttt{NULL}};
      \draw[->, thick, color=jet] (m1.east) -- (m2.west);
      \draw[->, thick, color=negative, dashed] (m2.east) -- (m3.west);
      \draw[->, thick, color=jet] (m3.east) -- (m4.west);

      \draw[->, thick, color=positive] (m1.south) -- (1.0, -0.9) -- (m3.south);
      \node[text=positive, below=0.5cm of m2] {bypass: \texttt{p->next = temp->next}};

      \node[text=neutral, above=0.15cm of m2] {doomed: \texttt{free(temp)}};
      \draw[->, thick, color=neutral] (2.0, 1.15) -- (m2.north);
    \end{tikzpicture}
  \end{center}
  \vspace{-0.2cm}
  \begin{itemize}
    \item Deleting 20: node 10's link jumps over it straight to 30 (solid
      green), then 20's box is freed. Skip the free and the node still sits
      in the heap with nobody pointing at it --- Week 1's leak, drawn.
  \end{itemize}
\end{frame}

\begin{frame}[fragile]{Search and Display: The Two Walks}
  \begin{block}{Both are one loop over \texttt{p = p->next}}
\begin{semiverbatim}
search(x):  p = head; while (p != NULL && p->data != x) p = p->next
            return (p != NULL)   // found or walked off the end
display():  p = head; while (p != NULL) { print p->data; p = p->next }
\end{semiverbatim}
  \end{block}
  \begin{itemize}
    \item Search is $O(n)$ worst case (the miss walks everything) and
      $O(1)$ best case (head holds \texttt{x}) --- name the case, per the
      Week 2 habit.
    \item Display never modifies: the traversal pointer \texttt{p} moves
      while \texttt{head} stays put. Moving \texttt{head} itself loses the
      list.
  \end{itemize}
\end{frame}

\transitionslide{Four operations banked. Now the classic: turn the chain around in place.}

% =============================================================================
% ACT 4 --- REVERSE + COSTS + CLOSE
% =============================================================================
\section{Reverse, Costs, Close}

\sectiondivider{Section 4}{Reverse, Costs, Close}

\begin{frame}{Reverse: Motivation Before Mechanism}
  \begin{itemize}
    \item The playlist just queued 500 songs oldest-first; the user taps
      ``play newest first.'' Copying all 500 into a second list doubles the
      memory for a rearrangement.
    \item The chain already holds every link we need --- each arrow just
      points the wrong way. Reversal re-points $n$ arrows using three
      pointers and $O(1)$ extra space.
    \item Socratic pause: with only \texttt{prev}, \texttt{curr}, and
      \texttt{next}, which arrow must be saved \emph{before} overwriting
      \texttt{curr->next} --- and why?
  \end{itemize}
  \begin{exampleblock}{The answer (check your instinct)}
    \texttt{next} $=$ \texttt{curr->next} first --- overwriting
    \texttt{curr->next} destroys the only path to the rest of the chain.
  \end{exampleblock}
\end{frame}

\begin{frame}[fragile]{Iterative Reverse: Three Pointers, One Pass}
  \begin{block}{The loop (Karumanchi Ch.~3, pp.74--162 PDF)}
\begin{semiverbatim}
prev = NULL; curr = head
while (curr != NULL):
    next = curr->next   // save the rest of the chain
    curr->next = prev   // flip this arrow
    prev = curr         // advance both walkers
    curr = next
head = prev
\end{semiverbatim}
  \end{block}
  \begin{itemize}
    \item Each node visited exactly once: $T(n) = n$ link flips --- $O(n)$
      worst case, $O(1)$ auxiliary space (three pointers, however long the
      list).
    \item Trace 10 $\to$ 20 $\to$ 30 once on paper before the lab: after the
      second iteration the chain reads 20 $\to$ 10 $\to$ \texttt{NULL} with
      \texttt{curr} parked at 30.
  \end{itemize}
\end{frame}

\begin{frame}{Reverse Mid-Pass, Drawn}
  \begin{center}
    \begin{tikzpicture}[every node/.style={font=\scriptsize}]
      \node[draw, fill=white, minimum width=0.9cm, minimum height=0.65cm,
        align=center] (r1) at (0.5, 0) {10};
      \node[draw, fill=white, minimum width=0.9cm, minimum height=0.65cm,
        align=center] (r2) at (2.2, 0) {20};
      \node[draw, fill=white, minimum width=0.9cm, minimum height=0.65cm,
        align=center] (r3) at (3.9, 0) {30};
      \node[draw, fill=neutral!20, minimum width=0.9cm, minimum height=0.65cm,
        align=center] (r4) at (5.6, 0) {\texttt{NULL}};
      \draw[<-, thick, color=positive] (r1.east) -- (r2.west);
      \draw[->, thick, color=jet] (r2.east) -- (r3.west);
      \draw[->, thick, color=jet] (r3.east) -- (r4.west);

      \node[text=neutral] at (0.5, -0.85) {\texttt{prev}};
      \draw[->, thick, color=neutral] (0.5, -0.7) -- (r1.south);
      \node[text=primary-blue] at (2.2, -0.85) {\texttt{curr}};
      \draw[->, thick, color=primary-blue] (2.2, -0.7) -- (r2.south);
      \node[text=primary-gold] at (3.9, -0.85) {\texttt{next}};
      \draw[->, thick, color=primary-gold] (3.9, -0.7) -- (r3.south);

      \node[text=positive, above=0.15cm of r1] {flipped};
      \node[text=neutral, above=0.15cm of r3] {still forward};
    \end{tikzpicture}
  \end{center}
  \vspace{-0.2cm}
  \begin{itemize}
    \item Snapshot after one flip: 10's arrow points back at \texttt{NULL},
      \texttt{prev} sits on 10, \texttt{curr} on 20, \texttt{next} saved 30.
      The frontier between flipped and forward walks right until
      \texttt{curr} $=$ \texttt{NULL}.
  \end{itemize}
\end{frame}

\begin{frame}{Singly Linked List: The Full Price List}
  \begin{center}
  {\footnotesize
  \begin{tabular}{lll}
    \toprule
    \textbf{operation} & \textbf{bound} & \textbf{note} \\
    \midrule
    insert at front & $O(1)$ worst case & two writes; order matters \\
    insert at end / position & $O(n)$ worst case & the walk dominates \\
    delete (value known by walk) & $O(n)$ worst case & bypass itself $O(1)$ \\
    search & $O(n)$ worst / $O(1)$ best & miss walks everything \\
    display & $\Theta(n)$ & must visit all $n$ \\
    iterative reverse & $O(n)$ worst, $O(1)$ space & one pass, three pointers \\
    \bottomrule
  \end{tabular}}
  \end{center}
  \begin{itemize}
    \item The house rule, Week 2 style: every bound above names its operation
      and its case. ``Insert is $O(1)$'' without saying \emph{where} is half
      an answer.
  \end{itemize}
\end{frame}

\begin{frame}{Three Classic Mistakes}
  \begin{itemize}
    \item \bad{Rewriting \texttt{head} before saving it} --- in front-insert
      (step order) and in reverse (unsaved \texttt{next}): the chain's only
      entry address is gone, and the heap still bills you for every orphaned
      node.
    \item \bad{Walking off the end} --- \texttt{p = p->next} past
      \texttt{NULL} dereferences nothing: check \texttt{p != NULL} in every
      loop condition, and test empty and single-node lists first.
    \item \bad{Freeing and still pointing} --- using \texttt{temp} after
      \texttt{free(temp)}, or forgetting the free entirely: Week 1's
      dangling pointer and leak, now wearing list clothing.
  \end{itemize}
  \begin{exampleblock}{How to revise}
    Re-trace the insert, delete, and reverse snapshots with a covered answer
    column --- the exam is traces, not definitions.
  \end{exampleblock}
\end{frame}

\begin{frame}{Week 6 Summary}
  \begin{itemize}
    \item \key{Chain over slots}: unpredictable size plus anywhere-edits beat
      $O(1)$ indexing --- that is when the list wins; otherwise the array
      keeps its crown.
    \item \key{Node}: \texttt{struct node \{ int data; struct node *next;
      \};} with \texttt{head} as sole entry and \texttt{NULL} as terminator;
      singly, doubly, and circular shapes for different jobs.
    \item \key{Five operations}: front-insert $O(1)$ worst case; end/position
      insert, delete, and search $O(n)$ worst case; display $\Theta(n)$;
      three-pointer reverse $O(n)$ worst case in $O(1)$ space.
  \end{itemize}
  \begin{exampleblock}{Next week}
    The backward link (\texttt{p->prev}), circular lists, the stack and queue
    reborn on chains --- and polynomials that only store what is nonzero.
  \end{exampleblock}
\end{frame}

\begin{frame}{Before You Go: Predict}
  \begin{block}{Three one-line predictions}
    Answer from the algorithms, then check in the lab.
  \end{block}
  \begin{enumerate}
    \item List \texttt{head} $\to$ 5 $\to$ 15 $\to$ \texttt{NULL}. Insert 10 at
      the front: write the two assignments in the safe order.
    \item Delete 15 above: which link is rewritten, and what is freed?
    \item Reverse \texttt{head} $\to$ 1 $\to$ 2 $\to$ 3 $\to$ \texttt{NULL}:
      after the loop, what does \texttt{head} point at, and what is the cost?
  \end{enumerate}
  \muted{Answers: \texttt{newNode->next = head}, then
    \texttt{head = newNode}. Rewrite 5's \texttt{next} to \texttt{NULL};
    free 15's node. \texttt{head} $\to$ 3 $\to$ 2 $\to$ 1 $\to$
    \texttt{NULL}; $O(n)$ worst case, $O(1)$ space.}
\end{frame}

\begin{frame}[allowframebreaks]{References}
  \footnotesize
  \bibliographystyle{plain}
  \bibliography{../../Bibliography_base}
\end{frame}

\end{document}
